% Firstpass — verified cascade routing.
%
% Build:  pdflatex main.tex && pdflatex main.tex
%
% EVERY empirical number in this paper is a \result{} placeholder until it is filled from a
% committed artifact under docs/benchmarks/. Unfilled placeholders render in red as [TODO: ...]
% so a draft can never be mistaken for a finished result, and \fillcheck at the end of the
% document reports how many remain. Do not replace a placeholder with a number you have not
% read out of a committed artifact.
\documentclass[11pt]{article}

\usepackage[margin=1in]{geometry}
\usepackage{amsmath,amssymb}
\usepackage{booktabs}
\usepackage{xcolor}
\usepackage{hyperref}
\usepackage{url}

\newcounter{todocount}
% Robust, because these appear inside \caption{}, which is a moving argument and would otherwise
% expand the counter step at the wrong time.
\DeclareRobustCommand{\result}[1]{\stepcounter{todocount}\textcolor{red}{\textbf{[TODO: #1]}}}
\newcommand{\fillcheck}{%
  \ifnum\value{todocount}>0
    \begin{center}\textcolor{red}{\textbf{DRAFT --- \arabic{todocount} unfilled results.
    Not submittable.}}\end{center}
  \fi}

\title{Verified Cascade Routing:\\
Serving the Cheapest Model That Provably Passes a Gate,\\
With an Auditable Record of Why}

\author{Shekhar Mudarapu\\ \texttt{firstpass}}
\date{}

\begin{document}
\maketitle

\begin{abstract}
Most LLM routing systems decide \emph{before} generation: a classifier reads the prompt and
predicts which model will suffice. We study the complementary regime --- deciding \emph{after}
generation, by running a verifier on the actual output and escalating only on failure --- and ask
what such a system must provide to be trustworthy in production rather than merely cheap.

The mechanism is not new: cheapest-first cascades with verification date to FrugalGPT, and AutoMix
established that self-verification is noisy enough to need a meta-verifier. Our contribution is
threefold and orthogonal to routing accuracy. First, a \textbf{distribution-free bound} on the rate
at which wrong answers are served, obtained by split-conformal calibration of the serve threshold
on held-out data. Second, a \textbf{tamper-evident audit record}: every routing decision, gate
verdict, and cost is hash-chained so an external auditor can re-derive the entire decision history
from stored records alone --- a property no system in a survey of ${\sim}130$ routing methods
provides. Third, \textbf{zero-retrain model onboarding}: adding a model is a configuration change,
where predictive routers require retraining a policy.

We evaluate on MBPP with real unit-test gates and report results under the metric of RouterBench,
the field's standard router benchmark. We also report a negative structural finding: RouterBench,
and benchmarks like it, \emph{cannot} evaluate verification routing at all, because they store one
response per model with no executable oracle. Finally, we quantify the regime in which our own
headline result holds by degrading the gate deliberately, and find that the serve threshold the production policy uses ($\tau = 1.0$) rejects roughly half of
all \emph{correct} cheap answers, and that a lower threshold is both cheaper and no less accurate.
\end{abstract}

\fillcheck

\section{Introduction}

An LLM serving stack must decide which model answers each request. The dominant framing is
\emph{predictive}: featurise the prompt, estimate which model will be adequate, dispatch. This is
the shape of RouteLLM, Hybrid LLM, and the large majority of the routing literature. It has a
structural weakness that is easy to state and hard to remove: the decision is made at the one moment
when the evidence that would settle it --- the answer --- does not yet exist. When the prediction is
wrong, the failure surfaces in production, and nothing in the system explains why that model was
chosen.

The alternative is to decide afterwards. Generate with the cheapest model, run a real check against
the output, and escalate only if the check fails. Where ``correct'' is machine-checkable --- code
with tests, structured output against a schema, extraction with a known answer --- this converts
routing from a prediction problem into a verification problem, and the cost saving becomes a
\emph{consequence} of the cheap model usually being sufficient rather than a bet that it will be.

We are not the first to make that observation, and we want to be precise about what is and is not
new here (\S\ref{sec:related}). What we claim is that a verification router becomes deployable only
once it answers three questions the literature has not addressed, because they are engineering
properties rather than accuracy properties:

\begin{enumerate}
  \item \textbf{How often does it serve a wrong answer?} Not empirically on the authors' data ---
  with a bound that holds on the operator's own traffic, under no distributional assumptions.
  \item \textbf{Can the decisions be audited?} If a routing decision is contested weeks later, can
  an auditor with only the stored records reconstruct what happened and detect tampering?
  \item \textbf{What does adding a model cost?} A predictive router retrains. A verification router
  should not have to.
\end{enumerate}

\paragraph{Contributions.}
\begin{itemize}
  \item A verified cascade router with a split-conformal served-failure bound
  (\S\ref{sec:method}), implemented as a production OpenAI-compatible proxy rather than a research
  prototype.
  \item A hash-chained audit record over routing decisions that an external party can re-derive
  (\S\ref{sec:audit}).
  \item An evaluation under RouterBench's own metric (AIQ over the non-decreasing convex hull),
  against its Zero Router bar (\S\ref{sec:results-aiq}).
  \item A negative result: standard router benchmarks are structurally unable to evaluate
  verification routing (\S\ref{sec:why-not-routerbench}).
  \item An imperfect-gate study that quantifies where our own guarantee stops holding
  (\S\ref{sec:results-sweep}), and an evaluation of AutoMix's meta-verifier in that regime
  (\S\ref{sec:results-meta}).
\end{itemize}

\section{Related work}
\label{sec:related}

We situate this against the taxonomy of Awesome-Routing-LLMs, which organises ${\sim}130$ routing
methods by \emph{when} the decision is made: pre-judgment, verification, memory-based, and
system-level analysis. Our work sits in the smallest of these families, verification routing by
self-assessment.

\begin{table}[h]
\centering
\small
\begin{tabular}{@{}p{0.24\textwidth}p{0.36\textwidth}p{0.32\textwidth}@{}}
\toprule
\textbf{Work} & \textbf{Mechanism} & \textbf{Difference here} \\
\midrule
FrugalGPT & Cheapest-first cascade; a learned scorer on the answer decides escalation. & Operator-authored gate (their tests, their schema); audit record; calibrated bound. \\
AutoMix & Few-shot self-verification plus a POMDP meta-verifier treating verification as noisy. & External deterministic gates where the domain allows; we evaluate their meta-verifier in our setting (\S\ref{sec:results-meta}). \\
CP-Router & Conformal prediction on output probabilities sizes a prediction set \emph{before} committing. & Same tool, different quantity: our bound is on served-failure \emph{after} the gate. \\
RouteLLM, Hybrid LLM & Prompt-conditioned classifier selects a model. & No prompt-conditioned prediction; adding a model requires no retraining. \\
RouterBench, RouterEval & Benchmarks for router comparison. & We adopt RouterBench's metric; \S\ref{sec:why-not-routerbench} shows why its dataset cannot score verification routers. \\
\bottomrule
\end{tabular}
\caption{Nearest prior art. The cascade-with-verification mechanism is established; our
contributions are the guarantee, the audit record, and the onboarding property.}
\end{table}

To the best of our knowledge, no method in that survey emits a tamper-evident record of its routing
decisions. The survey's safety category concerns \emph{attacks} on routers (adversarial redirection,
supply-chain compromise) rather than auditability of routing itself.

\section{Method}
\label{sec:method}

\subsection{Verified cascade}

Let $L = (m_1, \dots, m_k)$ be a ladder of models ordered by increasing price, and let
$g : \mathcal{Y} \rightarrow [0,1]$ be a gate assigning a score to an output. For request $x$, we
generate $y_1 = m_1(x)$ and serve it if $g(y_1) \geq \tau$; otherwise we generate $y_2 = m_2(x)$ and
repeat. Cost accrues for every rung attempted, including rejected ones.

The gate is supplied by the operator and is domain-specific: a unit-test suite, a JSON schema, an
LLM judge, or $k$-sample self-consistency. This is the point of departure from predictive routing:
$g$ reads $y$, not $x$.

\subsection{Calibrating the serve threshold}

The standing objection to any cascade is that $\tau$ is a hand-tuned hyperparameter with no
guarantee. We replace it with a calibrated one. Given held-out pairs $(g(y_i), c_i)$ where $c_i$
indicates whether $y_i$ was actually correct, we choose the lowest $\tau$ such that the empirical
failure rate among served items, inflated by a Hoeffding upper confidence bound, is at most
$\alpha$:
\[
\hat{r}(\tau) + \sqrt{\frac{\ln(1/\delta)}{2n_\tau}} \;\leq\; \alpha ,
\qquad
\hat{r}(\tau) = \frac{|\{i : g(y_i) \geq \tau,\ \neg c_i\}|}{n_\tau} .
\]
Serving on $g(y) \geq \tau$ then carries a served-failure rate of at most $\alpha$ with probability
at least $1 - \delta$. The bound is distribution-free: no Gaussian assumption, no assumption that
calibration and deployment traffic are identically distributed beyond exchangeability. It is
conservative, so it never under-covers.

Critically, calibration and evaluation must not share data. We calibrate on one split and report
realized served-failure on a disjoint split.

\subsection{Auditability}
\label{sec:audit}

Every request emits a record containing the ladder, the rung served, each gate verdict and score,
the cost, and the hash of the previous record. Records are chained over a canonical JSON encoding
that does not depend on struct field order or build features, so an auditor holding only the stored
records can recompute every hash and verify the linkage. Tampering with any historical decision
invalidates every subsequent hash.

This is the property we believe is genuinely absent from the literature. A predictive router's
decision is a forward pass whose inputs are not retained; there is nothing to re-derive. A verified
cascade's decision is a sequence of discrete, recorded events, which makes the audit record
possible at all.

\section{Why not RouterBench}
\label{sec:why-not-routerbench}

The natural evaluation is to run on RouterBench, the standard benchmark for router comparison. We
report that this is not possible, for reasons that generalise to router benchmarks as a class.

RouterBench provides, for each of ${\sim}36{,}500$ prompts, one stored response per model, a
quality score, and a cost. This is complete for a predictive router: read the prompt, choose a
model, look up its score. It is insufficient for a verification router in two independent ways.

\textbf{No executable oracle.} A verification router must run a real gate on the output. RouterBench
stores a quality score but no mechanism to derive one independently, so the only available ``gate''
is the stored score --- which is also the ground truth. Gate and oracle would be the same
measurement, and the evaluation would measure nothing.

\textbf{The code subset is not re-scorable.} MBPP is the one subset where an executable gate could
in principle be reconstructed from the canonical dataset's tests. RouterBench's MBPP prompts,
however, are paraphrases of the canonical items (mean similarity $\approx 0.78$ to the nearest
canonical text) and carry no test cases or function-name constraints, so stored responses define
arbitrary function names and cannot be run against canonical asserts. The \texttt{sample\_id} field
suggests an index mapping, but it is not consistent enough to validate.

We therefore invert the comparison: rather than evaluating our router on their data, we evaluate
\emph{their metric} on ours (\S\ref{sec:results-aiq}). We suggest that a benchmark able to score
verification routers must ship executable oracles with a visible/hidden split --- which, to our
knowledge, none currently does.

\section{Evaluation}
\label{sec:eval}

\paragraph{Setup.} MBPP, $n = 974$ tasks. Ladder: \texttt{claude-haiku-4-5} $\rightarrow$ \texttt{claude-opus-4-8}. The gate is the visible unit-test prefix; ground truth is a disjoint
hidden test set, so the gate can be wrong in both directions and we can measure it. Candidate code
executes in a network-isolated, fail-closed sandbox. Each task is solved once per rung and every
policy is replayed over that identical matrix, so policy comparisons carry no sampling noise and
are paired by construction.

\paragraph{Reproducibility.} The measured matrix is committed. Every table below is regenerated by
\texttt{firstpass-bench --replay <checkpoint>}, which requires no API key, no sandbox, and no
spend.

\subsection{Does routing beat picking one model?}
\label{sec:results-policy}

\begin{table}[h]
\centering
\begin{tabular}{@{}lrrrr@{}}
\toprule
\textbf{Policy} & \textbf{Success} & \textbf{\$/success} & \textbf{Escalated} & \textbf{Converted} \\
\midrule
always-cheap  & 0.78 & \$0.0018 & --- & --- \\
always-top    & 0.94 & \$0.0079 & --- & --- \\
first-pass    & 0.93 & \$0.0033 & 19\% & 78\% \\
\bottomrule
\end{tabular}
\caption{Paired over identical tasks. Paired $\Delta$ vs always-top: success
$-0.010$ (95\% CI $[-0.021, -0.001]$, 6 wins / 16 losses) and cost $-\$0.00440$ per task
(95\% CI $[-0.00469, -0.00412]$). Both intervals exclude zero: first-pass is significantly
cheaper \emph{and} significantly, if slightly, less accurate. The marginal intervals overlap and
would suggest otherwise --- which is why the paired test is the one reported.}
\end{table}

\subsection{Under RouterBench's metric}
\label{sec:results-aiq}

We compute AIQ, $\mathrm{AIQ}(R) = \frac{1}{c_{\max}-c_{\min}} \int_{c_{\min}}^{c_{\max}}
\tilde{q}_R(c)\, dc$, over the non-decreasing convex hull, exactly as defined by RouterBench, and
compare against their \emph{Zero Router}: the parameter-free interpolation between the raw models.
RouterBench reports that no learned router significantly outperformed this baseline, and that
routers were \emph{worse} than it on MBPP specifically --- making it a pre-registered bar rather
than one chosen after seeing our result.

\begin{table}[h]
\centering
\begin{tabular}{@{}lr@{}}
\toprule
\textbf{Curve} & \textbf{AIQ} \\
\midrule
Zero Router (models only)   & 0.8634 \\
{}+ first-pass              & 0.9161 \\
Oracle (perfect foresight)  & 0.9361 \\
\bottomrule
\end{tabular}
\caption{AIQ lift over the Zero Router: $+0.0526$, closing 72\% of the headroom to perfect
foresight. RouterBench reports no learned router significantly clearing this bar, and routers
performing \emph{worse} than it on MBPP specifically.}
\end{table}

\subsection{When the gate is allowed to be wrong}
\label{sec:results-sweep}

Our headline result holds under a near-perfect gate. To find its boundary we degrade the gate
deliberately: the production policy serves at $\tau = 1.0$ (all visible tests pass), and lowering
$\tau$ serves partial-credit answers, manufacturing genuine false accepts with ground truth
attached. This obtains real gate error without buying a weaker verifier.

\begin{table}[h]
\centering
\begin{tabular}{@{}rrrrrr@{}}
\toprule
$\tau$ & \textbf{Served cheap} & \textbf{False-accept} & \textbf{False-reject} & \textbf{Success} & \textbf{\$/success} \\
\midrule
1.00 & 59\% & 0.030 & \textbf{0.530} & 0.9271 & \$0.02136 \\
0.70 & 73\% & 0.034 & 0.323 & \textbf{0.9271} & \textbf{\$0.01872} \\
0.60 & 75\% & 0.036 & 0.276 & 0.9261 & \$0.01824 \\
0.50 & 82\% & 0.041 & 0.000 & 0.9230 & \$0.01697 \\
0.00 & 100\% & 0.209 & 0.000 & 0.7906 & \$0.01380 \\
\bottomrule
\end{tabular}
\caption{Judged run ($k{=}5$), the only configuration measured with a \emph{graded} gate score;
the unit-test gate alone is binary on MBPP. Gate AUC (score $\rightarrow$ oracle-correct): 0.938.
At $\tau = 1.0$ the gate rejects \textbf{53\%} of cheap answers that were in fact correct, buying
escalations that were never needed. $\tau = 0.70$ attains \emph{identical} success at 12\% lower
cost per success, so the production threshold is dominated.}
\end{table}

On the opus ladder the split-conformal bound is \emph{feasible} and holds out of sample:
$\lambda = 0.50$ calibrated on one half (empirical risk 0.0379, serving 81\%), with a held-out
realized served-failure of \textbf{0.0521} over 403 served items, against $\alpha = 0.10$.

RouterBench reports that cascading routers beat the Zero Router while judge error stays at or below
$0.1$, degrading sharply past $0.2$. Our measured gate sits inside that regime: false-accept is
0.030--0.041 across the thresholds that preserve success, an order of magnitude below the point at
which they observe cascades degrading.

\subsection{Does a meta-verifier help?}
\label{sec:results-meta}

AutoMix argues that verification is noisy enough that its output should be treated as evidence
rather than truth, and introduces a meta-verifier --- deliberately non-LLM, so verifier error does
not compound --- to decide under that uncertainty. We implement the same idea: bucket the evidence,
estimate $P(\text{correct} \mid \text{evidence})$ on a calibration split, and escalate when
$P < p^\ast$, where $p^\ast$ is derived from the ladder's measured price ratio rather than tuned.

\begin{table}[h]
\centering
\begin{tabular}{@{}lrrrr@{}}
\toprule
\textbf{Policy} & \textbf{Success} & \textbf{Total \$} & \textbf{Wins} & \textbf{Losses} \\
\midrule
first-pass     & 0.9343 & \$1.4674 & --- & --- \\
meta-verified  & 0.9343 & \$1.4674 & 0 & 0 \\
\bottomrule
\end{tabular}
\caption{Held-out comparison on the opus ladder ($p^\ast = 0.816$). The meta-verifier makes
\emph{identical} decisions to plain first-pass here --- 0 wins, 0 losses --- which is the expected
result where the gate is a near-perfect oracle: with a false-reject rate of 0.000 there is no
verifier error for a meta-verifier to correct. It is reported rather than omitted, because a
negative result on the workload where the mechanism cannot help is the honest scope of AutoMix's
contribution, not evidence against it.}
\end{table}

\section{Limitations}

\textbf{Verification needs something to verify against.} This approach is strongest where
correctness is checkable. For open-ended prose the available gates are an LLM judge or $k$-sample
self-consistency, and in our own measurements a judge gate did not pay for itself once its calls
were metered. If an output has no check, this is the wrong tool.

\textbf{One domain.} Our results are on MBPP. We do not claim they transfer to other task
distributions, and the field's own benchmarks cannot currently test that claim
(\S\ref{sec:why-not-routerbench}).

\textbf{The zero-regression result is workload-dependent.} Escalation can only cause harm when the
gate rejects a correct answer and the next rung then fails. On MBPP the first condition essentially
never occurred. \S\ref{sec:results-sweep} is the honest boundary of that claim, not a footnote to
it.

\textbf{Ladder sensitivity.} Cost savings depend on the price of the top rung, which is chosen by
the operator. We report savings against more than one ceiling for this reason.

\section{Conclusion}

Verification routing is a small and well-precedented corner of the routing literature. What has been
missing is not a better mechanism but the properties that make one deployable: a bound on how often
it is wrong, a record an auditor can check, and an onboarding cost that does not scale with the
model catalogue. We have implemented those and measured them under the field's own metric, and
reported both where the approach wins and the regime in which its strongest result stops holding.

\fillcheck

\begin{thebibliography}{9}
\bibitem{frugalgpt} L. Chen, M. Zaharia, J. Zou. FrugalGPT: How to Use Large Language Models While
Reducing Cost and Improving Performance. arXiv:2305.05176, 2023.
\bibitem{automix} P. Aggarwal, A. Madaan, et al. AutoMix: Automatically Mixing Language Models.
NeurIPS 2024, arXiv:2310.12963.
\bibitem{cprouter} J. Su, et al. CP-Router: An Uncertainty-Aware Router Between LLM and LRM.
AAAI 2025, arXiv:2505.19970.
\bibitem{routellm} I. Ong, et al. RouteLLM: Learning to Route LLMs with Preference Data.
arXiv:2406.18665, 2024.
\bibitem{hybridllm} D. Ding, et al. Hybrid LLM: Cost-Efficient and Quality-Aware Query Routing.
ICLR 2024, arXiv:2404.14618.
\bibitem{routerbench} Q. Hu, et al. RouterBench: A Benchmark for Multi-LLM Routing System.
arXiv:2403.12031, 2024.
\bibitem{mbpp} J. Austin, et al. Program Synthesis with Large Language Models. arXiv:2108.07732,
2021.
\end{thebibliography}

% Machine-readable draft state, so a build can be gated on it rather than on someone remembering
% to look: grep the build output for this line before submitting anywhere.
\typeout{FIRSTPASS_UNFILLED_RESULTS=\arabic{todocount}}

\end{document}
